\section{Adaptive Large Neighbourhood Search}
\label{sec:method}

Our method is an adaptive large neighbourhood search (ALNS) \cite{shaw1998using,ropke2006adaptive,pisinger2010large}:
it repeatedly removes a few tasks from the current plan and reinserts them, and it accepts the result by simulated
annealing \cite{kirkpatrick1983optimization}. Its strength comes from a plan representation in which every move is
feasible and cheap to evaluate exactly.

\subsection{Plan representation}
\label{sec:repr}

A plan is a pair $(\pi, C)$ of a global order $\pi$ of the tasks and a coalition $C_j$ for every task; every agent
visits its tasks in the order of $\pi$. A forward pass over $\pi$ computes the schedule in time
$O(\sum_j |C_j|)$: for the next task $j$, each member's arrival follows from its previous task, $s_j$ is their
maximum and $f_j = s_j + d_j$. The representation has three properties.
\begin{enumerate}
\item \emph{Deadlock-free.} Every precedence between two tasks on a route follows $\pi$, so no cycle can form.
\item \emph{Exact.} We compute travel times exactly as the benchmark's simulator does, so the forward pass reproduces
  the simulator's makespan bit for bit; our tests check this by replaying plans in the simulator, and every reported
  result is the simulator's own score.
\item \emph{Minimal covers.} Since the simulator does not let an agent join a task that is already covered, every
  $C_j$ must be a minimal cover: no member can be left out without breaking coverage. The search keeps this invariant.
\end{enumerate}
Restricting the search to such plans loses nothing:
\begin{proposition}
For every deadlock-free plan there is a plan $(\pi, C)$ with minimal covers whose makespan is not larger.
\end{proposition}
\begin{proof}[Proof sketch]
Let $\pi$ order the tasks by start time, breaking ties by a topological order of the route precedences (which exists
because the plan does not deadlock); every route follows $\pi$, and the forward pass gives the same schedule. Now drop
redundant members one at a time. Task $j$ then starts at the maximum of fewer arrivals, so not later. The dropped
member $i$ goes straight from its previous location $p$ to its next task $k$; it used to arrive at $k$ at
$f_j + \tau_{jk} \ge \alpha_{ij} + \tau_{jk} = f_p + \tau_{pj} + \tau_{jk} \ge f_p + \tau_{pk}$ by the triangle
inequality, so it arrives no later. Since starts are monotone in arrivals, no task starts later and no agent returns
later.
\end{proof}

\subsection{Exact insertion in linear time}
\label{sec:insertion}

The central operation inserts a task $j$ into a plan that lacks it: it chooses a position in $\pi$ and a coalition.
Let the position be after the prefix $\pi_{1..p}$. The prefix schedule does not change, so the arrival of every agent
at $j$ is known. We take the \emph{earliest-arrival cover}: add agents that contribute to an unmet requirement in
order of arrival until $j$ is covered, then drop redundant members, latest arrival first. This is a minimal cover
with the earliest possible start $s_j$ at that position.

The new makespan follows in $O(|C_j|)$. The earliest-start schedule is a longest-path computation in the acyclic graph
of the routes, so let the \emph{tail} $\theta_k$ of a task be the length of the longest path from its start to the
end: $\theta_k = d_k$ plus the largest, over its members, of the travel to that member's next task plus that task's
tail, or of the travel home. Inserting $j$ adds paths through $j$ and replaces each member's route arc $p \to k$ by
$p \to j \to k$, which is never shorter by the triangle inequality. Hence the makespan after the insertion is exactly
\begin{equation}
  M' = \max\Big(M,\; s_j + d_j + \max_{i \in C_j} \big(\tau_{j k_i} + \theta_{k_i}\big)\Big),
  \label{eq:insert}
\end{equation}
where $M$ is the old makespan, $k_i$ is member $i$'s next task after the position ($\tau_{jk_i} + \theta_{k_i}$ is
replaced by the travel home if there is none), and the tails are those of the old plan, since no task after the
position has a path to $j$. One sweep over $\pi$ updates the agents' arrivals at $j$ incrementally and scores every
position; positions where no capable agent's arrival changed are skipped, and for plans with more than 64 positions
only 64 are tried: the first, the last and those next to $j$'s nearest tasks.

The search objective is the makespan plus $\lambda$ times the mean finish time ($\lambda = 0.1$), which breaks the
many ties in makespan and rewards schedules with slack. For the second term, the delays that the new coalition causes
at its members' next tasks give a lower bound for every position; positions are then evaluated best-first by this
bound, with a forward pass that re-runs only agents whose state differs from the old schedule and stops as soon as
none does, until the bound exceeds the best position found. The chosen insertion is thus exact, and the whole
evaluation is identical to a full forward pass (which our tests check).

\subsection{Search}
\label{sec:search}

\begin{algorithm}[t]
\caption{ALNS worker (one process)}
\label{alg:alns}
\begin{algorithmic}[1]
\STATE $x \gets$ best of randomized constructions, run for 5\% of the budget
\STATE $x^* \gets x$; operator weights $w \gets 1$; temperature $t \gets t_0$
\WHILE{time remains}
  \STATE pick destroy $D$ and repair $R$ by roulette on $w$; $q \sim U\{1, \dots, q_{\max}\}$
  \STATE $x' \gets R(D(x, q))$ \COMMENT{or exchange of route tails}
  \IF{$g(x') \le g(x)$ or $u < e^{-(g(x') - g(x))/t}$}
    \STATE $x \gets x'$; with probability 0.1 re-sort $\pi$ by start time
    \IF{$x$ is better than $x^*$ (makespan, then $g$)} \STATE $x^* \gets x$ \ENDIF
  \ENDIF
  \STATE update scores; every 100 iterations update $w$; cool $t$
  \STATE after 3000 iterations without a new best: $x \gets x^*$
  \STATE every second: exchange $x^*$ with the other workers
\ENDWHILE
\RETURN $x^*$
\end{algorithmic}
\end{algorithm}

Algorithm~\ref{alg:alns} shows one worker; $g$ is the search objective.

\paragraph{Construction.} For the first 5\% of the budget, a worker builds plans by list scheduling
\cite{kolisch1996serial} (repeatedly append the task whose earliest-arrival cover starts first), by the same with noisy scores, and, up to 100 tasks, by
random-order best insertion, and keeps the best.

\paragraph{Destroy operators} remove $q$ tasks, with $q$ uniform in $\{1, \dots, q_{\max}\}$ and
$q_{\max} = \mathrm{clip}(\mathrm{round}(4\sqrt{50/m}), 2, 6)$: large instances are short of iterations and gain from
many small moves. \emph{Random} removes uniformly; \emph{related} removes tasks close to a random seed task in a random
mix of distance, start time and requirement \cite{shaw1998using}; \emph{worst wait} prefers tasks whose members wait
longest; \emph{critical} removes tasks of the chain that determines the makespan (following the latest-arriving member
back to its depot) and their spatial neighbours; \emph{route segment} removes consecutive tasks of one agent's route,
half of the time the agent that returns last. Selection is randomised by rank as in \citet{ropke2006adaptive}. A sixth
move exchanges the route tails of two agents with identical traits after a random position in $\pi$; it keeps all
covers minimal and needs no repair.

\paragraph{Repair operators} reinsert the removed tasks one at a time with the exact insertion of
Section~\ref{sec:insertion}: in random order, in random order with arrival times perturbed by up to 20\% when
choosing covers, largest requirement first, or by regret-2 (for up to ten removed tasks, else in random order).

\paragraph{Acceptance and adaptation.} Candidates are accepted by simulated annealing on $g$, with the temperature
cooled geometrically over the wall-clock budget from $t_0 = 0.01\, g_0 \min(1, (50/m)^{1.5})$ to $t_0 / 50$, where
$g_0$ is the objective of the construction. Operator weights adapt by segment scores \cite{ropke2006adaptive}. After
an accepted move the order $\pi$ is re-sorted by start time with probability 0.1; this is the same plan and schedule,
but later insertions then see positions on a time line. After 3000 iterations without a new best plan, the search
returns to the best plan.

\paragraph{Parallel workers.} On $c$ cores, $c$ independent workers with different seeds exchange their best plans
through shared memory every second: a worker publishes its best plan if it beats the shared one and adopts the shared
plan if it is strictly better. The kernels are compiled with Numba; a worker runs \ItsFiftylo--\ItsFiftyhi{} iterations per
second on the 50-task instances and about \ItsFivehundred{} on the 500-task ones. All parameters were set on a separate tuning split
(Section~\ref{sec:setup}); several other components (elite pools with crossover, kicks, growing removal sizes, member
swaps, delay-aware covers) were tried there and not kept.
